\documentclass[12pt, a4paper]{article}

\usepackage{setspace} % Necessary for setting space
\usepackage{amsmath}
\usepackage{amssymb}
\usepackage{harvard} % Necessary for Econometrica style reference
\usepackage{hyperref} % Necessary for citing websites
\usepackage{lasstyle}
\usepackage{booktabs} % Table
\usepackage{graphicx} % Table

\newcounter{headercounter}
\setcounter{headercounter}{1}

\newcommand{\header}[1]{\begin{center}{\large \textbf{\Roman{headercounter}. #1}} \end{center} \stepcounter{headercounter}}

% \renewcommand{\thesection}{\Roman{section}} % Roman section enumeration

\doublespacing

\author{Lucas Sernik}
\title{Spending pattern differences between mayors who might be reelected and those who are binded by term limits: A case study from municipalities in São Paulo}
\date{\today}

\begin{document}
\maketitle

\begin{abstract}
This paper examines whether reelection incentives distort municipal spending composition in São Paulo, Brazil. Using a panel of 643 municipalities over 2017--2024 and exploiting variation in incumbent reelection eligibility generated by Brazil's two-term limit rule, I estimate a time fixed effects model relating reelection eligibility to the share of expenditure allocated to six functional categories. I find limited support for the political budget cycle hypothesis: reelection eligibility is significant only for administrative spending, where term-limited mayors allocate approximately one percentage point more to administration, consistent with reduced electoral discipline over bureaucratic overhead. No significant effects are detected for health, education, or other service categories. The discretionary share of spending is negative and significant for health, education, and administration, suggesting that earmarked federal transfers are the primary drivers of spending composition, leaving limited scope for politically motivated budget manipulation. These findings highlight the constraining role of fiscal institutions on political budget cycles.
\end{abstract}

\newpage

\header{Introduction}

In a democratic society, citizens vote in a set of candidates for government positions who are supposed to make decisions that promote the public good. Recent developments in Political Science and Political Economy offer different views on whether the existence of term limits provide any societal benefit. Smart and Sturn (2013) develop a game that models voters and self-interested politicians, where term limits help voters screen out politicians who use their position for their own self interest, improving overall welfare despite allowing for sub-optimal behavior on the binding term. Fowler (2024) discusses how term limits might have unintended effects when applied to political positions. He argues that term limits might remove from office a high-quality official that voters like and who voters would prefer to stay in office and, more importantly, that the mere existence of a term limit modifies the incentives that the official have to do their job. Such concern is tested by Besley and Case (1995). 

They make a compelling study that provides empirical evidence to the vast theoretical literature on politicians and what drives their decisions. They argue that the existence of term limits might distort politicians incentives concerning which policies are implemented. After setting their theoretical framework and some of its shortcomings, they analyze state governors through the years 1950 to 1986 and conclude that ``lame ducks'', governors who cannot run for the next term, increase their spending and taxations in their last term, which might be impopular, since they generally do not need to mantain their political reputation. They point out that term limits might have adverse effects on economic policy, arguing that politicians might behave differently and be more inefficient due to the different incentives that such rule imposes on their constrained maximization problem.

This paper will focus on Brazil, more specifically in municipalities in the state of São Paulo. Brazil transitioned from military rule to a democratic regime at the end of the 1980's, with a movement for direct elections in 1985 and the creation of a new constitution in 1988. Recent efforts for greater transparency in government actions and budgeting led Brazilian congress to pass the ``Lei da Transparência'' in 2009 and the ``Lei de Acesso à Informação'' in 2011, which obliges public institutions to publish detailed documents with information on government revenues and expenses on the internet. In the state of São Paulo, expenditure data are available dating back to 2014. This paper will attempt to exploit expense data made available by these laws in the 645 municipalities of the state of São Paulo from years 2017 to 2024. This time period is the largest one that contains entire election terms (two) up to 2026. I will employ similar methods to Besley and Case (1995) to test whether mayors binded by term limits behave differently than those who may be reelected in the next period, according to the type of expense and the year. Preliminary results suggest that some categories like education and health see little difference in spending, since most of its funding come from earmarked federal and state funds. However, in funds where mayors can change spending patterns, in administration specifically.

\header{Data}

Expense data was collected through the publicly available API from Tribunal de Contas do Estado de São Paulo \cite{TCESP:expenses}. Out of the 645 municipalities in the state of São Paulo, the public API offers expense data for 644, all municipalities except the São Paulo municipality, which has their own website and data available from 2021 onwards, reason why it will be dropped from the dataset. Out of the 644 remaining municipalities in the state of São Paulo, the public API offers expense data from 2014 until 2026. Municipal elections in Brazil happen every four years since 1996 \cite{TSE:eleicao}. This means that municipal elections happend in 2012, 2016, 2020 and 2024. I restricted calls to years 2017 to 2024, the two complete election periods with relevant data. After downloading with the API, the municipality of Florínea was dropped, because it did not contain expense for the entire time period. This left us with 643 municipalities, which is the number of municipalities in the analysis.

Election data was collected through the publicly available data from Tribunal Superior Eleitoral \cite{TSE:data}. Through their data transparency tool, I manually collected data for the years 2016 and 2020, the years when elections occured. I first filtered for the mayor position and municipalities in the state of São Paulo. Then, I selected for municipality, reelection status, counting situation (whether the candidate was elected or not), and party symbol. On metrics, chose quantity of candidates and type of file csv (en-US). The file contained all the candidates for all municipalities in a given election year. It was then processed further to have a column that identified whether the winner of the 2016 was reelected, if the incumbent from 2016 ran in 2020, if the 2020 winner was reelected, and total number of candidates in years 2016 and 2020 for each municipality.

Control data was collected from 2 sources. The first is from Instituto Brasileiro de Geografia e Estatística \cite{IBGE:sidra}. I was able to make API calls and collect population, GPD per capita, basic sanitation services and age distribution at the municipal granularity from the year 2010, the most recent data available. The second is from Atlas Brasil \cite{atlas:idh}. They compile data from other Brazilian institutions and offer to the public. It was taken manually from their website after filtering for municipalities in the state of São Paulo and choosing the most recent available data, also from 2010.

\header{Estimation and Methods}

Let $s_{ikt}$ denote the share of total expenditure allocated to category $k$ by municipality $i$ in year $t$. The baseline model is:

\begin{equation}
    s_{ikt} = \beta_0 
            + \beta_1 R_i 
            + \beta_2 D_{ikt} 
            + \beta_3 (R_i \times D_{ikt}) 
            + \beta_4 E_{it} 
            + \mathbf{X}_i'\boldsymbol{\gamma} 
            + \mathbf{P}_i'\boldsymbol{\delta} 
            + \mu_t 
            + \varepsilon_{it}
\end{equation}

\noindent where:

\begin{itemize}
    \item $R_i \in \{0,1\}$ is an indicator equal to 1 if the incumbent mayor elected in 2016 in municipality $i$ is eligible for reelection in 2020, and 0 otherwise
    \item $D_{ikt} \in [0,1]$ is the share of expenditure in category $k$, in municipality $i$ and year $t$ funded by non-earmarked sources (\textit{tesouro} and \textit{recursos próprios}), measuring fiscal discretion
    \item $R_i \times D_{ikt}$ is the interaction term, capturing whether reelection incentives have a stronger effect on spending composition in a given category when the mayor has greater fiscal discretion
    \item $E_{it}$ is total expenditure per capita in municipality $i$ and year $t$
    \item $\mathbf{X}_i$ is a vector of pre-determined socioeconomic controls measured at the 2010 Census, including PIB per capita, municipal HDI, sanitation access, and age distribution shares
    \item $\mathbf{P}_i$ is a vector of political controls, including the number of candidates in the 2016 election and an indicator for whether the 2016 winner was running as an incumbent
    \item $\mu_t$ denotes year fixed effects
    \item $\varepsilon_{it}$ is the idiosyncratic error term
\end{itemize}

\noindent The coefficient of interest is $\beta_1$, which identifies the effect of reelection eligibility on spending composition. The interaction coefficient $\beta_3$ tests whether this effect is moderated by fiscal discretion.

This is the model that will be discussed in the results. I am also exploring the Random Effects model and a difference in difference model where cities in which mayors cannot be reelected are used as control.

\newpage

\header{Results}

\begin{table}[htbp]
\centering
\caption{Effect of Reelection Eligibility on Municipal Spending Shares}
\label{tab:fe_results}
\resizebox{\textwidth}{!}{%
\begin{tabular}{lcccccc}
\hline\hline
 & \multicolumn{6}{c}{Dependent Variable: Spending Share (\%)} \\
\cmidrule(lr){2-7}
 & Health & Education & Urbanismo & Social Assist. & Administration & Transport \\
\hline
$R_i$ & 0.332 & $-$0.027 & 0.308 & $-$0.001 & $-$0.982$^{**}$ & 0.198 \\
 & (0.347) & (0.402) & (0.358) & (0.149) & (0.398) & (0.222) \\
$D_{ikt}$ & $-$0.022$^{*}$ & $-$0.142$^{***}$ & $-$0.014 & 0.009 & $-$0.141$^{***}$ & $-$0.006 \\
 & (0.012) & (0.030) & (0.010) & (0.006) & (0.038) & (0.004) \\
$R_i \times D_{ikt}$ & $-$0.011 & 0.020 & $-$0.006 & $-$0.021 & 0.115$^{**}$ & 0.006 \\
 & (0.026) & (0.040) & (0.019) & (0.027) & (0.054) & (0.006) \\
Expenditure p.c. & $-$0.000$^{***}$ & $-$0.000$^{***}$ & 0.000$^{*}$ & 0.000$^{***}$ & 0.000$^{***}$ & 0.000$^{***}$ \\
 & (0.000) & (0.000) & (0.000) & (0.000) & (0.000) & (0.000) \\
PIB p.c. & 0.000 & $-$0.000 & $-$0.000$^{**}$ & $-$0.000$^{**}$ & $-$0.000$^{**}$ & 0.000$^{*}$ \\
 & (0.000) & (0.000) & (0.000) & (0.000) & (0.000) & (0.000) \\
IDHM & $-$10.808$^{*}$ & 11.925$^{*}$ & 4.790 & $-$11.856$^{***}$ & $-$25.333$^{***}$ & $-$29.334$^{***}$ \\
 & (5.591) & (6.697) & (5.556) & (2.051) & (7.010) & (3.366) \\
\% Children & $-$0.128$^{***}$ & 0.574$^{***}$ & 0.074 & $-$0.056$^{***}$ & $-$0.265$^{***}$ & $-$0.119$^{***}$ \\
 & (0.048) & (0.055) & (0.046) & (0.016) & (0.059) & (0.032) \\
N Candidates 2016 & $-$0.236$^{**}$ & 0.103 & 0.004 & $-$0.156$^{***}$ & $-$0.095 & $-$0.035 \\
 & (0.097) & (0.116) & (0.091) & (0.027) & (0.099) & (0.057) \\
\hline
Observations & 5,144 & 5,144 & 5,116 & 5,144 & 5,144 & 3,726 \\
Municipalities & 643 & 643 & 641 & 643 & 643 & 524 \\
R$^2$ (Overall) & 0.071 & 0.317 & 0.024 & 0.177 & 0.069 & 0.138 \\
Year FE & Yes & Yes & Yes & Yes & Yes & Yes \\
\hline\hline
\multicolumn{7}{l}{\footnotesize Clustered standard errors at municipality level in parentheses.} \\
\multicolumn{7}{l}{\footnotesize $^{*}p<0.10$, $^{**}p<0.05$, $^{***}p<0.01$.} \\
\multicolumn{7}{l}{\footnotesize $R_i$: reelection eligibility indicator. $D_{ikt}$: demeaned discretionary} \\
\multicolumn{7}{l}{\footnotesize spending share by category. Standard errors clustered at municipality level.} \\
\end{tabular}%
}
\end{table}

\textbf{{\large Main Findings}}

Table \ref{tab:fe_results} presents the estimates from the time fixed effects panel model for six spending categories: health, education, urbanismo, social assistance, administration, and transport. The coefficient of interest, $\beta_1$, captures the effect of reelection eligibility on the share of municipal expenditure allocated to each category, conditional on fiscal discretion, municipal characteristics, and year fixed effects.

Across the six categories examined, reelection eligibility is statistically significant only for administrative spending ($\beta_1 = -0.98$, $p = 0.014$). Municipalities governed by reelection-eligible mayors allocate approximately one percentage point less of their total expenditure to administration relative to municipalities governed by term-limited mayors. This finding is consistent with the political budget cycle hypothesis of \cite{rogoff:1990}, which predicts that incumbents facing reelection shift resources away from less visible government consumption toward more electorally salient spending. However, no corresponding significant increase is detected in health, education, urbanismo, social assistance, or transport, suggesting that the reallocation away from administration does not translate into a clear increase in any single visible spending category. The null results for health and education are particularly notable given that these are the categories most associated with electoral returns in the Brazilian municipal context.

\textbf{\large{The Role of Fiscal Discretion}}

The coefficient on the demeaned discretionary share, $\beta_2$, is negative and statistically significant for health ($\beta_2 = -0.022$, $p = 0.063$), education ($\beta_2 = -0.142$, $p < 0.001$), and administration ($\beta_2 = -0.141$, $p < 0.001$). The consistent negative sign indicates that categories funded to a greater extent by discretionary treasury resources receive a lower share of total expenditure. This result reflects the institutional structure of Brazilian municipal finance: health and education spending is largely driven by constitutionally earmarked federal transfers --- specifically FUNDEB and SUS --- rather than by discretionary treasury allocations. Municipalities that rely more heavily on treasury funds tend to spread discretionary resources across a broader range of functions, mechanically reducing the share allocated to any individual category.

The interaction term $\beta_3$, which captures whether the effect of reelection eligibility on spending composition is moderated by fiscal discretion, is statistically significant only for administration ($\beta_3 = 0.115$, $p = 0.033$). The positive sign indicates that the negative effect of reelection eligibility on administrative spending is attenuated in municipalities with greater fiscal discretion. When mayors have more room to allocate treasury funds freely, the electoral incentive to reduce administrative overhead is weaker --- potentially because discretionary administrative spending serves political functions, such as patronage, that are not easily compressed even under electoral pressure. For all other categories the interaction term is small and statistically indistinguishable from zero, suggesting that fiscal discretion does not moderate the relationship between electoral incentives and spending in service delivery categories.

\textbf{{\large Control Variables}}

Several control variables yield findings of independent interest. The municipal HDI is negative and significant for administration ($\beta = -25.33$, $p < 0.001$), social assistance ($\beta = -11.86$, $p < 0.001$), and transport ($\beta = -29.33$, $p < 0.001$), indicating that wealthier municipalities allocate proportionally less to these categories --- consistent with economies of scale in administration and lower demand for social assistance in more developed municipalities. The share of children in the population is positive and significant for education ($\beta = 0.574$, $p < 0.001$) and negative and significant for administration ($\beta = -0.265$, $p < 0.001$), social assistance ($\beta = -0.056$, $p < 0.001$), and transport ($\beta = -0.119$, $p < 0.001$), which validates the demographic controls and confirms that spending composition responds to population needs. The number of candidates in the 2016 election --- a proxy for political competition --- is negative and significant for social assistance ($\beta = -0.156$, $p < 0.001$) and health ($\beta = -0.236$, $p = 0.015$), suggesting that more competitive political environments are associated with lower spending on redistributive categories, consistent with wealthier and more urbanized municipalities tending to have both more candidates and lower social assistance needs.

\textbf{{\large Model Specification}}

The F-test for poolability rejects the null of a common intercept across time periods in all regressions ($p < 0.001$), confirming that year fixed effects are appropriate and that aggregate time trends --- including changes in federal transfer rules and macroeconomic conditions --- meaningfully affect spending composition. Standard errors are clustered at the municipality level throughout to account for serial correlation in spending decisions within municipalities over time.

\textbf{{\large Discussion}}

Taken together, the results provide limited support for the political budget cycle hypothesis in the context of São Paulo municipalities. The only category for which reelection eligibility predicts spending composition is administration, where term-limited mayors spend more --- consistent with reduced electoral discipline over bureaucratic overhead. The absence of significant effects in health, education, and other service categories suggests that the institutional constraints imposed by earmarked federal transfers substantially limit the margin within which electoral incentives can distort spending. This finding contributes to the growing literature on fiscal federalism and political accountability in decentralized systems, highlighting that the effectiveness of electoral discipline over public spending depends critically on the degree of fiscal discretion available to local governments. In the Brazilian context, where constitutional earmarking covers the majority of municipal revenue in the two largest spending categories, the scope for politically motivated budget manipulation appears to be largely confined to administrative and overhead spending rather than frontline service delivery.

\bibliographystyle{econometrica}
\bibliography{Lucas_Sernik_third_draft}
\end{document}